% ============================================================================
%  Preamble of the thesis sample (pdfLaTeX).  Layout follows Phuc's template:
%  A4, 14pt, margins 35/20/30/30 mm, page number centred at the top.
% ============================================================================
\usepackage[utf8]{vietnam}
\usepackage[left=35mm,right=20mm,top=30mm,bottom=30mm,headsep=10mm]{geometry}
\usepackage{amsmath,amssymb,amsthm,mathtools}
\usepackage{times}
\usepackage[varg]{newtxmath}
\usepackage{bm}
\usepackage{microtype}
\usepackage{setspace}
\onehalfspacing
\usepackage{indentfirst}
\setlength{\parindent}{1cm}
\setlength{\parskip}{3pt plus 1pt}
\setlength{\emergencystretch}{1.5em}

\usepackage{graphicx}
\graphicspath{{figures/}}
\usepackage[font=small,labelfont=bf,labelsep=period,justification=justified,singlelinecheck=true]{caption}
\usepackage{subcaption}
\usepackage{booktabs,multirow,array,tabularx}
\usepackage[section,above]{placeins}
\usepackage[output-decimal-marker={,},group-separator={.},detect-all]{siunitx}
\usepackage[shortlabels]{enumitem}
\setlist{itemsep=2pt,topsep=4pt,parsep=0pt}
\usepackage{xcolor}
\usepackage{tikz}
\usetikzlibrary{positioning,arrows.meta,calc,fit,backgrounds,decorations.pathreplacing}
\usepackage[most]{tcolorbox}
\usepackage{algorithm}
\usepackage[noend]{algpseudocode}
\usepackage{titlesec}
\usepackage[titles]{tocloft}
\usepackage{fancyhdr}
\usepackage[hidelinks]{hyperref}
\usepackage[nameinlink,noabbrev]{cleveref}

% ---------------------------------------------------------------- colours
\definecolor{planframe}{HTML}{2a78d6}
\definecolor{planback}{HTML}{f3f7fd}
\definecolor{inktwo}{HTML}{52514e}

% ---------------------------------------------------------------- headings
\renewcommand{\chaptername}{CHƯƠNG}
\titleformat{\chapter}[display]
  {\centering\bfseries\fontsize{16}{20}\selectfont}
  {CHƯƠNG \thechapter}{4pt}{\MakeUppercase}
\titlespacing*{\chapter}{0pt}{0pt}{18pt}
\titleformat{\section}[hang]{\bfseries\fontsize{14.4}{18}\selectfont}{\thesection.}{0.4em}{}
\titlespacing*{\section}{0pt}{14pt}{6pt}
\titleformat{\subsection}[hang]{\bfseries\itshape\fontsize{14.4}{18}\selectfont}{\thesubsection.}{0.4em}{}
\titlespacing*{\subsection}{0pt}{10pt}{4pt}
\setcounter{secnumdepth}{2}
\setcounter{tocdepth}{2}

% un-numbered chapters (Mở đầu, Kết luận, ...) centred in capitals
\newcommand{\unchapter}[1]{%
  \chapter*{#1}\addcontentsline{toc}{chapter}{#1}\markboth{#1}{#1}}

% ---------------------------------------------------------------- page style
\pagestyle{fancy}
\fancyhf{}
\fancyhead[C]{\small\thepage}
\renewcommand{\headrulewidth}{0pt}
\fancypagestyle{plain}{\fancyhf{}\fancyhead[C]{\small\thepage}\renewcommand{\headrulewidth}{0pt}}

% ---------------------------------------------------------------- lists (TOC, LOF, LOT)
\AtBeginDocument{%
  \renewcommand{\contentsname}{MỤC LỤC}%
  \renewcommand{\listfigurename}{DANH MỤC CÁC HÌNH}%
  \renewcommand{\listtablename}{DANH MỤC CÁC BẢNG}%
  \renewcommand{\figurename}{Hình}%
  \renewcommand{\tablename}{Bảng}%
  \renewcommand{\chaptername}{CHƯƠNG}%
  \renewcommand{\bibname}{TÀI LIỆU THAM KHẢO}%
  \renewcommand{\proofname}{Chứng minh}}
\renewcommand{\cftchappresnum}{Chương }
\settowidth{\cftchapnumwidth}{\bfseries Chương 3.\ }
\renewcommand{\cftchapaftersnum}{.}
\renewcommand{\cftsecaftersnum}{.}
\renewcommand{\cftsubsecaftersnum}{.}
\renewcommand{\cftfigpresnum}{Hình }
\renewcommand{\cftfigaftersnum}{.}
\settowidth{\cftfignumwidth}{Hình 3.10.\ }
\renewcommand{\cfttabpresnum}{Bảng }
\renewcommand{\cfttabaftersnum}{.}
\settowidth{\cfttabnumwidth}{Bảng 3.10.\ }
\setlength{\cftbeforechapskip}{6pt}

% ---------------------------------------------------------------- numbering
\numberwithin{equation}{chapter}
\numberwithin{figure}{chapter}
\numberwithin{table}{chapter}
\renewcommand{\theequation}{\thechapter.\arabic{equation}}

% ---------------------------------------------------------------- theorems
\newtheoremstyle{vnplain}{6pt}{6pt}{\itshape}{}{\bfseries}{.}{0.5em}{}
\newtheoremstyle{vndef}{6pt}{6pt}{\normalfont}{}{\bfseries}{.}{0.5em}{}
\newtheoremstyle{vnrem}{6pt}{6pt}{\normalfont}{}{\bfseries\itshape}{.}{0.5em}{}
\theoremstyle{vnplain}
\newtheorem{theorem}{Định lý}[chapter]
\newtheorem{lemma}[theorem]{Bổ đề}
\newtheorem{proposition}[theorem]{Mệnh đề}
\newtheorem{corollary}[theorem]{Hệ quả}
\theoremstyle{vndef}
\newtheorem{definition}[theorem]{Định nghĩa}
\newtheorem{example}[theorem]{Ví dụ}
\theoremstyle{vnrem}
\newtheorem{remark}[theorem]{Nhận xét}
\renewcommand{\proofname}{Chứng minh}
\renewcommand{\qedsymbol}{$\square$}

\crefname{theorem}{Định lý}{Định lý}
\crefname{lemma}{Bổ đề}{Bổ đề}
\crefname{proposition}{Mệnh đề}{Mệnh đề}
\crefname{corollary}{Hệ quả}{Hệ quả}
\crefname{definition}{Định nghĩa}{Định nghĩa}
\crefname{example}{Ví dụ}{Ví dụ}
\crefname{remark}{Nhận xét}{Nhận xét}
\crefname{figure}{Hình}{Hình}
\crefname{table}{Bảng}{Bảng}
\crefname{chapter}{Chương}{Chương}
\crefname{section}{Mục}{Mục}
\crefname{subsection}{Mục}{Mục}
\crefname{equation}{}{}
\crefname{algorithm}{Thuật toán}{Thuật toán}
\creflabelformat{equation}{(#2#1#3)}
% \Cref would upper-case the first letter itself (loops on T5 "Đ"): give explicit names
\Crefname{theorem}{Định lý}{Định lý}
\Crefname{lemma}{Bổ đề}{Bổ đề}
\Crefname{proposition}{Mệnh đề}{Mệnh đề}
\Crefname{corollary}{Hệ quả}{Hệ quả}
\Crefname{definition}{Định nghĩa}{Định nghĩa}
\Crefname{example}{Ví dụ}{Ví dụ}
\Crefname{remark}{Nhận xét}{Nhận xét}
\Crefname{figure}{Hình}{Hình}
\Crefname{table}{Bảng}{Bảng}
\Crefname{chapter}{Chương}{Chương}
\Crefname{section}{Mục}{Mục}
\Crefname{subsection}{Mục}{Mục}
\Crefname{equation}{}{}
\Crefname{algorithm}{Thuật toán}{Thuật toán}
\newcommand{\crefpairconjunction}{ và }
\newcommand{\crefmiddleconjunction}{, }
\newcommand{\creflastconjunction}{ và }
\newcommand{\crefrangeconjunction}{--}


% ---------------------------------------------------------------- algorithms
\floatname{algorithm}{Thuật toán}
\renewcommand{\thealgorithm}{\thechapter.\arabic{algorithm}}
\algrenewcommand\algorithmicrequire{\textbf{Đầu vào:}}
\algrenewcommand\algorithmicensure{\textbf{Đầu ra:}}
\algrenewcommand\algorithmicfor{\textbf{với}}
\algrenewcommand\algorithmicdo{\textbf{thực hiện}}
\algrenewcommand\algorithmicif{\textbf{nếu}}
\algrenewcommand\algorithmicthen{\textbf{thì}}
\algrenewcommand\algorithmicelse{\textbf{ngược lại}}
\algrenewcommand\algorithmicreturn{\textbf{trả về}}

% ---------------------------------------------------------------- "planned content" box
\newtcolorbox{dukien}[1][]{enhanced,breakable,colback=planback,colframe=planframe,
  boxrule=0.6pt,arc=2pt,left=7pt,right=7pt,top=5pt,bottom=5pt,
  fonttitle=\bfseries\small,coltitle=white,title={Nội dung dự kiến#1},
  fontupper=\small\setstretch{1.15}}
\newcommand{\trang}[1]{\hfill{\normalfont\itshape(khoảng #1 trang)}}

% ---------------------------------------------------------------- notation
\newcommand{\R}{\mathbb{R}}
\newcommand{\N}{\mathbb{N}}
\newcommand{\B}{\mathcal{B}}
\newcommand{\Om}{\Omega}
\newcommand{\T}{\top}
\newcommand{\pos}[1]{\left(#1\right)_{+}}
\newcommand{\sgn}{\operatorname{sign}}
\DeclareMathOperator{\intr}{int}
\DeclareMathOperator{\conv}{conv}
\DeclareMathOperator*{\argmin}{arg\,min}
\DeclareMathOperator*{\argmax}{arg\,max}
\newcommand{\bphi}{\bm{\varphi}}
\newcommand{\norm}[1]{\left\|#1\right\|}

% ---------------------------------------------------------------- code listings
\usepackage{listings}
\lstset{language=Python,basicstyle=\ttfamily\footnotesize\setstretch{1.0},
  keywordstyle=\bfseries\color{planframe},commentstyle=\itshape\color{inktwo},
  stringstyle=\color{inktwo},frame=single,rulecolor=\color{black!25},framesep=4pt,
  xleftmargin=6pt,xrightmargin=6pt,showstringspaces=false,columns=fullflexible,
  keepspaces=true,backgroundcolor=\color{planback},aboveskip=6pt,belowskip=6pt}
